% ID: SEL-D0021-OA-E01
% ESTADO: validada. S022, revisada con los criterios oficiales.
\textbf{a)} Al multiplicar filas por columnas se obtiene
$A^2=\begin{pmatrix}8&4\\1&5\end{pmatrix}$. Como $B$ es diagonal,
$B^3=\operatorname{diag}(-27,1)$. Al sumar,
\[
\boxed{A^2+B^3=\begin{pmatrix}-19&4\\1&6\end{pmatrix}}.
\]
\textbf{b)} La segunda fila de $A+B=\begin{pmatrix}-1&4\\1&0\end{pmatrix}$
hace que la primera fila de $X$ sea la segunda de $A-B$, es decir,
$(1,-2)$. Si $r$ es la segunda fila de $X$, la primera ecuación
por filas da $-(1,-2)+4r=(5,4)$. Por tanto
\[
\boxed{X=\begin{pmatrix}1&-2\\3/2&1/2\end{pmatrix}}.
\]
El producto recupera las filas $(5,4)$ y $(1,-2)$ de $A-B$.
